%% toc: P2 14 14.5,14.6,14.8
\section{Proofs for finite-input decision procedures and the arbitrary-language scope}
\label{sec:annex-p02-w02-p02-017-finite-input-scope}

Throughout, \(P\) is a set of atoms, tagged as in
\cref{def:p02-fbc-language}, and for \(Q\subseteq P\) the language
\(\operatorname{Form}(Q)\) is identified with the set of
\(\alpha\in\operatorname{Form}(P)\) satisfying
\(\operatorname{Var}(\alpha)\subseteq Q\).  Valuations take values in
\(\{0,1\}\).

\emph{Computability facts used.}  By \cref{thm:p04-tm-partial-recursive}
the computable functions of \cref{def:partial-computable} are exactly
the total partial recursive functions of \cref{def:p04-partial-recursive}.
That class is closed under composition and primitive recursion, and it
contains every primitive recursive function, as
the theorem of \cref{sec:turing-machines-02} also states.  The following
functions are primitive recursive, by the displayed recursions from
zero, successor, projections, addition and multiplication
(the example in \cref{sec:turing-machines-02}):
\[
\begin{gathered}
 \operatorname{pd}(0)=0,\quad \operatorname{pd}(n+1)=n;\qquad
 a\mathbin{\dot{-}}0=a,\quad a\mathbin{\dot{-}}(b+1)=\operatorname{pd}(a\mathbin{\dot{-}} b);\\
 \overline{\operatorname{sg}}(0)=1,\quad
 \overline{\operatorname{sg}}(n+1)=0;\qquad
 \operatorname{eq}(a,b)=\overline{\operatorname{sg}}
   \bigl((a\mathbin{\dot{-}} b)+(b\mathbin{\dot{-}} a)\bigr);\qquad
 \max(a,b)=a+(b\mathbin{\dot{-}} a).
\end{gathered}
\]
Thus \(\operatorname{eq}(a,b)=1\) iff \(a=b\), and otherwise
\(\operatorname{eq}(a,b)=0\).  A machine that transforms its input string
and halts may be followed by a second such machine: rename the second
machine's states apart, return the head to cell \(0\), and identify the
first machine's accepting state with the second machine's start state.
The encoding paragraph of \cref{sec:turing-machines-01} supplies halting machines for writing
\(\operatorname{bin}(m)\) and for reading a binary word as a number, and
\cref{def:p02-sat-taut} supplies the encoding of a formula string as a
terminating linear-time transformation.  Writing a fixed list of
characters before and after a given string is likewise a halting
transformation.

\emph{Size of code numbers.}  If \(\alpha\in\operatorname{Form}(P_{\N})\) contains the
variable \(p_m\), then the encoding of \(\alpha\) contains the binary
digits of \(m\), so the binary numeral of \(\ulcorner\alpha\urcorner\) has
more than \(\abs{\operatorname{bin}(m)}\) digits.  A number whose numeral
has \(L\) digits is at least \(2^{L-1}\), and
\(m<2^{\abs{\operatorname{bin}(m)}}\) for every \(m\), including
\(m=0\).  Hence
\begin{equation}
 \ulcorner\alpha\urcorner>m
 \qquad\text{whenever }p_m\text{ occurs in }\alpha .
 \tag{S1}
\end{equation}

\emph{A decidable witness relation for \(K\).}  The two proofs of
\cref{sec:p02-w02-p02-017-infinite-premises} share the following
construction.  By \cref{thm:universal-computation} the universal function, written
\(U(e,x)\) there and not to be confused with the output function
\(\mathsf U\) below, is partial computable, hence partial recursive by
\cref{thm:p04-tm-partial-recursive}.  The function
\(F(x)\simeq U(x,x)\simeq\varphi_x(x)\) is its composition with two
projections, hence partial recursive; here \(\varphi_x\) is the
partial function with index \(x\) in the program numbering of
\cref{sec:universal-reductions-01},
not a formula.  The normal form
\cref{thm:p04-kleene-normal-form} with \(k=1\) gives an index \(e^{*}\)
such that
\[
 F(x)\simeq\mathsf U\bigl(\mu y\,\mathsf T_1(e^{*},x,y)\bigr)
 \qquad(x\in\N).
\]
Write \(T(x,y)\) for \(\mathsf T_1(e^{*},x,y)\).  Its characteristic
function is the composition of the primitive recursive characteristic
function of \(\mathsf T_1\) with a constant and projections, so \(T\) is
decidable.  The minimization is defined exactly when some \(y\) satisfies
\(T(x,y)\), and it then returns the least such \(y\).  Therefore
\begin{equation}
 x\in K
 \iff \varphi_x(x)\text{ is defined}
 \iff F(x)\text{ is defined}
 \iff \exists y\,T(x,y).
 \tag{S2}
\end{equation}

\begin{annexitem}{annex:p02-w02-p02-017-finite-reduction}{thm:p02-w02-p02-017-finite-reduction}{Finite inputs over an arbitrary set of atoms}
\begin{proof}
\emph{Clause 1.}  Induct on \(\alpha\).  For an atom \(q\),
\(\operatorname{Var}(q)=\{q\}\) and \(\widehat v(q)=v(q)=w(q)=\widehat
w(q)\).  For \(\bottom\) both values are \(0\).  Since
\(\operatorname{Var}(\neg\alpha)=\operatorname{Var}(\alpha)\) and
\(\operatorname{Var}(\alpha\circ\beta)=\operatorname{Var}(\alpha)\cup
\operatorname{Var}(\beta)\), agreement on the larger set gives agreement
on each immediate subformula; the induction hypotheses and the common
truth clause give equal values.

\emph{Clause 2.}  The valuation \(s^{0}\) is defined by cases on the
membership \(q\in A\); no choice is involved.  If some valuation \(v\)
satisfies \(\Gamma_0\), let \(s\) be the restriction of \(v\) to \(A\).
Then \(s^{0}\) agrees
with \(v\) on \(A\supseteq\operatorname{Var}(\gamma)\) for every
\(\gamma\in\Gamma_0\), so by clause~1 \(s^{0}\) satisfies \(\Gamma_0\).
The converse is immediate.  In the same way, if a valuation \(v\)
satisfies \(\Gamma_0\) and gives \(\varphi\) value \(0\), then
\(s^{0}\) does the same for that restriction \(s\), because
\(\operatorname{Var}(\varphi)\subseteq A\); so
\(\Gamma_0\not\models_P\varphi\) iff some \(s^{0}\) is a
countervaluation.  For finite \(A\), induction on \(\abs A\) shows that
there are exactly \(2^{\abs A}\) functions \(A\to\{0,1\}\); when
\(A=\varnothing\) the only one is the empty function and \(s^{0}\) is
the constant \(0\) valuation.

\emph{Clause 3.}  Induct on \(\alpha\in\operatorname{Form}(Q)\).  For an
atom \(q\in Q\), \(\widehat u(q^{c})=u(p_{c(q)})=v(q)=\widehat v(q)\).
For \(\bottom\) both values are \(0\).  The translation commutes with
\(\neg,\land,\lor,\to\), and both recursive truth definitions use the
same clauses, so the induction hypotheses carry over.  Given \(v\) on
\(P\), define \(u(p_n)=v(q)\) when \(n=c(q)\) for some \(q\in Q\), and
\(u(p_n)=0\) otherwise; injectivity of \(c\) makes the \(q\) unique, so
\(u\) is a single defined function.  Given \(u\) on \(P_{\N}\), define
\(v(q)=u(p_{c(q)})\) for \(q\in Q\) and \(v(q)=0\) for
\(q\in P\setminus Q\).  In both directions the partners satisfy the
hypothesis.  If \(v\) satisfies \(\Delta\), its partner \(u\) satisfies
\(\Delta^{c}\); if \(u\) satisfies \(\Delta^{c}\), its partner \(v\)
satisfies \(\Delta\).  For \(\Delta\cup\{\psi\}\subseteq
\operatorname{Form}(Q)\), a countervaluation to \(\Delta\models_P\psi\)
yields one to \(\Delta^{c}\models_{P_{\N}}\psi^{c}\) and conversely, by the same
partners.

\emph{Clause 4.}  Since \(A\) is finite there is a bijection
\(A\to\{0,\ldots,n-1\}\), which is an atom coding of \(A\).  Fix any atom
coding \(c\) of \(A\).  By clause~3 with \(Q=A\), \(\Gamma_0\) is
satisfiable iff \(\Gamma_0^{c}\) is, and \(\Gamma_0\models_P\varphi\) iff
\(\Gamma_0^{c}\models_{P_{\N}}\varphi^{c}\).  The algorithm of
\cref{thm:p02-decidability} decides finite entailment for its encoded
inputs, and a finite set is satisfiable iff it does not entail
\(\bottom\).  Its answers on the encoded input over \(c\) are therefore
the truth values of statements about \(\Gamma_0\) and \(\varphi\) in
which \(c\) does not occur, so they coincide for all atom codings of
\(A\).  None of the clauses uses a choice principle or compactness.
\end{proof}
\end{annexitem}

\begin{annexitem}{annex:p02-w02-p02-017-no-coding}{thm:p02-w02-p02-017-no-coding}{Only countable well-orderable atom sets can be coded}
\begin{proof}
Let \(c:P\to\N\) be injective.  Put \(a\prec b\) iff \(c(a)<c(b)\).  This
is a strict linear order because \(<\) is one and \(c\) is injective.  If
\(S\subseteq P\) is nonempty, \(c[S]\) has a least element \(m\), and the
unique \(a\in S\) with \(c(a)=m\) is \(\prec\)-least in \(S\).  So
\(\prec\) well-orders \(P\), and \(c\) is the required injection into
\(\N\).

Let \(C:\operatorname{Form}(P)\to\{0,1\}^{*}\) be injective.  The map
\(q\mapsto q\), sending an atom to the atomic formula, is injective, and
so is \(\nu:\{0,1\}^{*}\to\N\) sending \(w\) to the number whose binary
numeral is \(1w\): the numeral of a positive number has no leading zero,
and deleting its first digit recovers \(w\).  The composite
\(q\mapsto\nu(C(q))\) is an atom coding, and the first paragraph applies.

If \(g:\mathcal P(\N)\to\N\) were injective, the map sending \(n\) to
\(g^{-1}(n)\) for \(n\) in the range of \(g\), and to \(\varnothing\)
otherwise, would be a surjection \(\N\to\mathcal P(\N)\), contrary to
\cref{thm:p01-cantor}.  So \(\mathcal P(\N)\) admits no injection into
\(\N\).

Finally, an algorithm reads finite strings over a finite input alphabet
\(\Sigma\).  Replacing each character by a fixed-width binary block, as in
\cref{def:p02-sat-taut}, injects \(\Sigma^{*}\) into \(\{0,1\}^{*}\).  An
algorithm with \(\operatorname{Form}(P)\) as its input domain would
require an injective representation of formulas by strings in
\(\Sigma^{*}\), hence an injection
\(\operatorname{Form}(P)\to\{0,1\}^{*}\), and then \(P\) would be
well-orderable and inject into \(\N\).
\end{proof}
\end{annexitem}

\begin{annexitem}{annex:p02-w02-p02-017-infinite-undecidable}{thm:p02-w02-p02-017-infinite-undecidable}{Satisfiability of decidable infinite premise sets is undecidable}
\begin{proof}
Use the relation \(T\) of (S2).  For
\(y\in\N\) put
\[
 \theta_y=\neg p_0\land(p_{y+1}\lor\neg p_{y+1}),
 \qquad
 \Gamma_x=\{p_0\}\cup\{\theta_y:T(x,y)\}.
\]

\emph{Uniform decidability.}  The function \(t(y)=\ulcorner\theta_y\urcorner\)
is computable: a halting machine writes \(\operatorname{bin}(y+1)\),
surrounds it with the fixed characters of the formula string of
\(\theta_y\), applies the encoding of \cref{def:p02-sat-taut}, and
prefixes the digit \(1\); the result is the numeral of \(t(y)\).  By
(S1), \(t(y)>y+1>y\).  Define by primitive
recursion
\[
 B(x,n,0)=0,\qquad
 B(x,n,y+1)=\max\bigl(B(x,n,y),\
   \operatorname{eq}(n,t(y))\cdot\operatorname{char}_T(x,y)\bigr).
\]
Then \(B(x,n,m)=1\) iff some \(y<m\) has \(n=t(y)\) and \(T(x,y)\).  Put
\[
 M(x,n)=\max\bigl(\operatorname{eq}(n,\ulcorner p_0\urcorner),\
   B(x,n,n+1)\bigr),
\]
a computable function by the closure facts.  Since
\(\alpha\mapsto\ulcorner\alpha\urcorner\) is injective, \(n\) is the code
of a member of \(\Gamma_x\) iff \(n=\ulcorner p_0\urcorner\) or
\(n=t(y)\) for some \(y\) with \(T(x,y)\); in the latter case
\(y<t(y)=n<n+1\), so the bounded search finds \(y\).  Hence \(M\) is the
characteristic function of
\(\{(x,\ulcorner\gamma\urcorner):\gamma\in\Gamma_x\}\).  For fixed \(x\),
\(n\mapsto M(x,n)\) is computable, so each \(\Gamma_x\) is decidable.

\emph{Satisfiability.}  If \(x\notin K\), no \(y\) satisfies \(T(x,y)\)
by (S2), so \(\Gamma_x=\{p_0\}\), which
the valuation with every value \(1\) satisfies.  If \(x\in K\), fix a
\(y\) with \(T(x,y)\).  Then \(p_0,\theta_y\in\Gamma_x\), and no valuation
makes both \(p_0\) and \(\neg p_0\) true, so \(\Gamma_x\) is
unsatisfiable.

\emph{Consequences.}  Thus \(\{x:\Gamma_x\text{ is satisfiable}\}=
\N\setminus K\).  By the corollary following
\cref{thm:halting-undecidable} in \cref{sec:universal-reductions-03}, this
set is not computably enumerable.  If it were decidable, the characteristic function
of \(K\), which is \(x\mapsto1\mathbin{\dot{-}}\operatorname{char}_{\N\setminus
K}(x)\), would be computable, contrary to
\cref{thm:halting-undecidable}.  A finite subset of \(\Gamma_x\) is a
finite subset of \(\operatorname{Form}(P_{\N})\), so it has a finite encoded input over the
identity coding of its variables, on which \cref{thm:p02-decidability}
decides satisfiability.
\end{proof}
\end{annexitem}

\begin{annexitem}{annex:p02-w02-p02-017-no-computable-valuation}{thm:p02-w02-p02-017-no-computable-valuation}{A decidable satisfiable premise set with no computable model}
\begin{proof}
Use \(F\), \(e^{*}\) and \(T\) of (S2).

\emph{Unique witnesses.}  Define by primitive recursion
\[
 N(x,0)=1,\qquad
 N(x,y+1)=N(x,y)\cdot\bigl(1\mathbin{\dot{-}}\operatorname{char}_T(x,y)\bigr),
\]
so \(N(x,y)=1\) iff no \(z<y\) satisfies \(T(x,z)\).  Let \(W(x,y)\) hold
iff \(T(x,y)\) and \(N(x,y)=1\); its characteristic function
\(\operatorname{char}_T\cdot N\) is computable.  If \(W(x,y)\) and
\(W(x,y')\) with \(y<y'\), then \(N(x,y')=1\) contradicts \(T(x,y)\); so
each \(x\) has at most one \(y\) with \(W(x,y)\), namely the least \(y\)
with \(T(x,y)\) when there is one.  By the normal form,
\begin{equation}
 W(x,y)\iff F(x)\text{ is defined and }y=\mu z\,T(x,z),
 \qquad\text{and then } F(x)=\mathsf U(y).
 \tag{S3}
\end{equation}

\emph{The premise set.}  Put \(\tau_y=p_{2y+1}\lor\neg p_{2y+1}\) and
\[
 \delta_{x,y}=
 \begin{cases}
  \neg p_{2x}\land\tau_y, & \mathsf U(y)=0,\\
  p_{2x}\land\tau_y, & \mathsf U(y)\ne0,
 \end{cases}
 \qquad
 \Delta=\{\delta_{x,y}:W(x,y)\}.
\]

\emph{Decidability.}  The functions
\(d_0(x,y)=\ulcorner\neg p_{2x}\land\tau_y\urcorner\) and
\(d_1(x,y)=\ulcorner p_{2x}\land\tau_y\urcorner\) are computable in the
same way as \(t\) in the preceding proof, and \(\mathsf U\) is primitive
recursive.  Hence
\[
 d(x,y)=\operatorname{eq}(\mathsf U(y),0)\cdot d_0(x,y)
 +\bigl(1\mathbin{\dot{-}}\operatorname{eq}(\mathsf U(y),0)\bigr)\cdot d_1(x,y)
 =\ulcorner\delta_{x,y}\urcorner
\]
is computable.  Since \(p_{2x}\) and \(p_{2y+1}\) occur in
\(\delta_{x,y}\), (S1) gives
\(d(x,y)>2x\ge x\) and \(d(x,y)>2y+1>y\).  Define by primitive recursion
\[
\begin{gathered}
 C(n,x,0)=0,\qquad
 C(n,x,y+1)=\max\bigl(C(n,x,y),\
   \operatorname{char}_W(x,y)\cdot\operatorname{eq}(n,d(x,y))\bigr),\\
 D(n,0)=0,\qquad
 D(n,x+1)=\max\bigl(D(n,x),\ C(n,x,n+1)\bigr).
\end{gathered}
\]
Then \(D(n,n+1)=1\) iff some \(x,y\le n\) satisfy \(W(x,y)\) and
\(n=d(x,y)\).  By injectivity of code numbers, \(n\) codes a member of
\(\Delta\) iff \(n=d(x,y)\) for some \((x,y)\) with \(W(x,y)\), and then
\(x,y<n\).  So \(n\mapsto D(n,n+1)\) is the characteristic function of
\(\{\ulcorner\delta\urcorner:\delta\in\Delta\}\), and \(\Delta\) is
decidable.

\emph{Satisfiability without choice.}  Define the valuation \(v\) on
\(P_{\N}\) by \(v(p_{2x})=1\) iff \(x\in K\) and \(F(x)\ne0\), and
\(v(p_{2y+1})=0\).  It is given by an explicit definition, so no choice
principle is used.  Every \(\tau_y\) is true under every valuation.  Let
\(\delta_{x,y}\in\Delta\), so \(W(x,y)\); by
(S3), \(x\in K\) and
\(F(x)=\mathsf U(y)\).  If \(\mathsf U(y)=0\), then \(F(x)=0\), so
\(v(p_{2x})=0\) and \(\neg p_{2x}\land\tau_y\) is true.  If
\(\mathsf U(y)\ne0\), then \(F(x)\ne0\), so \(v(p_{2x})=1\) and
\(p_{2x}\land\tau_y\) is true.  Hence \(v\) satisfies \(\Delta\).

\emph{No computable model.}  Suppose a valuation \(w\) satisfies
\(\Delta\) and \(n\mapsto w(p_n)\) is computable.  Then
\(h(x)=1\mathbin{\dot{-}} w(p_{2x})\) is computable and total.  By
\cref{thm:enumeration}(b) there is an index \(b\) with
\(\varphi_b\simeq h\).  Then \(\varphi_b(b)=h(b)\) is defined, so
\(F(b)=h(b)\), \(b\in K\), and by (S3)
there is a \(y\) with \(W(b,y)\) and \(\mathsf U(y)=h(b)\).  So
\(\delta_{b,y}\in\Delta\).  If \(h(b)=0\), then \(\delta_{b,y}=\neg
p_{2b}\land\tau_y\), so \(w(p_{2b})=0\) and \(h(b)=1\).  If \(h(b)=1\),
then \(\mathsf U(y)\ne0\), \(\delta_{b,y}=p_{2b}\land\tau_y\), so
\(w(p_{2b})=1\) and \(h(b)=0\).  Both cases are contradictory, so no
computable valuation satisfies \(\Delta\).
\end{proof}
\end{annexitem}

\begin{annexitem}{annex:p02-w02-p02-017-bpi-boundary}{thm:p02-w02-p02-017-bpi-boundary}{Finite checks supply the hypothesis, not the valuation}
\begin{proof}
\emph{Clause 1.}  Apply
\cref{thm:p02-w02-p02-017-finite-reduction}(2) to each finite
\(\Gamma_0\subseteq\Gamma\) with \(\varphi=\bottom\), so that
\(A=\operatorname{Var}(\Gamma_0)\).  The check succeeds on \(\Gamma_0\)
exactly when \(\Gamma_0\) is satisfiable.

\emph{Clause 2.}  Let \(c\) be an atom coding of \(P\).  The map
\(\alpha\mapsto\alpha^{c}\) is injective on \(\operatorname{Form}(P)\),
by induction on formulas: atoms go to distinct variables because \(c\) is
injective, and the outermost constructor and immediate subformulas of
\(\alpha\) are recovered from those of \(\alpha^{c}\).  Let \(\Gamma\) be
finitely satisfiable and let \(\Sigma\subseteq\Gamma^{c}\) be finite.
The set \(\Gamma_0=\{\gamma\in\Gamma:\gamma^{c}\in\Sigma\}\) is finite,
since translation maps it injectively into \(\Sigma\), and
\(\Gamma_0^{c}=\Sigma\).  \(\Gamma_0\) is satisfiable, so \(\Sigma\) is
satisfiable by \cref{thm:p02-w02-p02-017-finite-reduction}(3) with
\(Q=P\).  Thus \(\Gamma^{c}\subseteq\operatorname{Form}(P_{\N})\) is finitely satisfiable,
and \cref{thm:prop-compactness}(2) gives a valuation \(u\) satisfying
\(\Gamma^{c}\).  Its partner \(v\) on \(P\) from the same clause~3
satisfies \(\Gamma\).  \Cref{thm:prop-compactness} is proved for the
fixed alphabet \(P_{\N}\) by an explicit enumeration and uses no choice
principle, as recorded in \cref{thm:p01-bpi-propositional-compactness}
and at the opening of \cref{sec:p02-free-boolean-completeness}; the
translation steps are explicit definitions.

\emph{Clause 3, from \(\mathsf{BPI}\).}  Over
\(\mathsf{ZF}+\mathsf{BPI}\), \cref{thm:p02-fbc-completeness} states for
every set \(P\) that every finitely satisfiable set of formulas of
\(\operatorname{Form}(P)\) is satisfiable.

\emph{Clause 3, to \(\mathsf{BPI}\).}  Assume the scheme.  Let \(V\) be a
set and use the language of \cref{def:p01-bpi-propositional-compactness},
with variables \(p_v\) (\(v\in V\)), constants \(\bot,\top\) and
connectives \(\neg,\land,\lor\).  Take \(P=\{p_v:v\in V\}\) as a set of
atoms and translate each formula \(\sigma\) of that language to
\(\sigma^{\bullet}\in\operatorname{Form}(P)\) by
\(p_v^{\bullet}=p_v\), \(\bot^{\bullet}=\bottom\),
\(\top^{\bullet}=\neg\bottom\), and by commuting with
\(\neg,\land,\lor\).  For \(\nu:V\to\mathbf2\) let
\(v_\nu(p_v)=\nu(v)\).  Induction on \(\sigma\) shows that the value of
\(\sigma\) under \(\nu\) equals \(\widehat{v_\nu}(\sigma^{\bullet})\); the
case \(\top\) uses \(\widehat{v_\nu}(\neg\bottom)=1\).  Every valuation
on \(P\) is \(v_\nu\) for exactly one \(\nu\).  Let \(\Sigma\) be finitely
satisfiable in the sense of \cref{def:p01-bpi-propositional-compactness},
and let \(\Phi\subseteq\Sigma^{\bullet}\) be finite.  Each member of
\(\Phi\) has a nonempty set of preimages in \(\Sigma\); a finite family of
nonempty sets has a choice function in \(\mathsf{ZF}\), by induction on
the number of members.  The chosen preimages form a finite
\(\Sigma_0\subseteq\Sigma\), satisfied by some \(\nu\), and then
\(v_\nu\) satisfies \(\Phi\).  So \(\Sigma^{\bullet}\) is finitely
satisfiable, the scheme gives a valuation \(v_\nu\) satisfying it, and
\(\nu\) satisfies \(\Sigma\).  This is the compactness principle of
\cref{def:p01-bpi-propositional-compactness} for every set \(V\), and
\cref{thm:p01-bpi-propositional-compactness} yields \(\mathsf{BPI}\).

\emph{Consequence.}  Suppose \(\mathsf{ZF}\) proved that, for every set
\(P\) and every \(\Gamma\subseteq\operatorname{Form}(P)\), success of all
the finite checks on finite subsets of \(\Gamma\) implies that
\(\Gamma\) is satisfiable.  By clause~1 \(\mathsf{ZF}\) would prove the
scheme of clause~3, and hence \(\mathsf{BPI}\).  In the proof of
\cref{thm:p02-fbc-separation} the countervaluation is
\(v_U=\chi_U\eta_P\), where \(U\) is an ultrafilter extending a
prescribed proper filter; its existence is the single use of
\(\mathsf{BPI}\), through \cref{thm:p01-bpi-equivalence}, and no step of
that proof applies a decision procedure.  When \(P\) has no atom coding,
\cref{thm:p02-w02-p02-017-no-coding} shows that no algorithm even reads
the formulas whose classes form \(L_P\).
\end{proof}
\end{annexitem}

\begin{annexitem}{annex:p02-w02-p02-017-pairs}{thm:p02-w02-p02-017-pairs}{Valuations of the pair-selection premises are choice functions}
\begin{proof}
Each \(a\in P\) lies in exactly one member of \(\mathcal A\), by
disjointness; write \(A_a\) for it.  Since every member of
\(\mathcal A\) has two elements, the formulas of \(\Gamma_{\mathcal A}\)
are exactly \((a\lor b)\land\neg(a\land b)\) with \(\{a,b\}\in\mathcal
A\) and \(a\ne b\).  Such a formula is true under \(v\) iff
\(v(a)+v(b)=1\), that is, iff exactly one of \(a,b\) has value \(1\).
This proves the characterization of satisfying valuations.

\emph{Finite subsets.}  Let \(\Gamma_0\subseteq\Gamma_{\mathcal A}\) be
finite.  The map \(a\mapsto A_a\) sends the finite set
\(\operatorname{Var}(\Gamma_0)\) onto
\(\mathcal A_0=\{A\in\mathcal A:A\cap\operatorname{Var}(\Gamma_0)\neq
\varnothing\}\), so \(\mathcal A_0\) is finite.  A finite family of
nonempty sets has a choice function in \(\mathsf{ZF}\): the empty family
has the empty function, and a choice function on \(n\) members extends
to \(n+1\) members by one further choice.  Let \(g\) be a choice
function on \(\mathcal A_0\), and put \(v(a)=1\) iff \(A_a\in\mathcal A_0\)
and \(a=g(A_a)\).  Every formula of \(\Gamma_0\) has the form displayed
above with \(\{a,b\}\in\mathcal A_0\), and exactly one of \(a,b\), namely
\(g(\{a,b\})\), has value \(1\).  So \(v\) satisfies \(\Gamma_0\).

\emph{The bijection.}  For a satisfying valuation \(v\), let \(g_v(A)\)
be the unique \(a\in A\) with \(v(a)=1\); this defines a choice function
on \(\mathcal A\) by separation, without any choice.  For a choice
function \(g\) on \(\mathcal A\), put \(v_g(a)=1\) iff \(a=g(A_a)\); it
gives value \(1\) to exactly one element of each member, so it satisfies
\(\Gamma_{\mathcal A}\).  Then \(g_{v_g}=g\), and \(v_{g_v}=v\) because
\(v(a)=1\) iff \(a\) is the element of \(A_a\) with value \(1\).

\emph{Under \(\mathsf{BPI}\).}  The set \(\Gamma_{\mathcal A}\) is
finitely satisfiable, so over \(\mathsf{ZF}+\mathsf{BPI}\) it is
satisfiable by \cref{thm:p02-fbc-completeness}, and the bijection turns a
satisfying valuation into a choice function on \(\mathcal A\).  The case
\(\mathcal A=\varnothing\) is included: then \(P\) and
\(\Gamma_{\mathcal A}\) are empty, and the empty valuation corresponds to
the empty choice function.
\end{proof}
\end{annexitem}
